Short-range infrared links are attractive where radio-frequency operation is
constrained, but simple embedded receivers typically decode them with
hand-tuned amplitude or timing thresholds that discard a symbol, and with it an
entire byte, whenever a pulse falls outside its acceptance window. This paper
presents LORAD, an experimental optical proximity link, and shows that
replacing threshold decoding with maximum-likelihood sequence detection
substantially improves link reliability without any change to the transmitter,
the receiver electronics, or the modulation.

The link uses baseband on-off keying in which symbol identity is carried by
pulse width rather than amplitude, a choice validated by measurement: with the
receiver's transimpedance stage saturated, peak amplitude held to within 0.17\%
across every pulse recorded while pulse-width classes remained separated by
12--40 pooled standard deviations. A dual-photodiode
receiver streams per-pulse observations to a decoder that combines a learned
emission model, inter-pulse gap timing, and the protocol's own byte-framing and
printable-ASCII constraints in a hidden Markov model decoded by the Viterbi
algorithm.

The system is characterized over five data rates from 28.6 to 171.4~bit/s using
597 labeled transmissions, with a further 157 transmissions decoded in a live
closed loop on hardware. Held out by message, the sequence decoder raises the
transmission success rate from 72.2\% to 100\% at 28.6~bit/s and from 36.1\% to 69.4\% at
114.3~bit/s. In live operation at 114.3~bit/s it raises transmission success from
5.9\% to 44.9\% and halves character error rate, recovering 41\% of the
transmissions the threshold decoder lost. Analysis shows the dominant failure
mode is undetected pulses rather than misclassified ones, which per-symbol
classifiers cannot address by construction, and that every such loss leaves a
recoverable signature in the inter-pulse gap.
